\documentclass[10pt,a4paper]{amsart}
\usepackage[margin=19mm]{geometry}
\usepackage{amsmath,amssymb,mathtools}
\usepackage[scr=rsfs,cal=euler]{mathalfa}
\usepackage{xfrac}
\usepackage[unicode,hidelinks,bookmarks=false]{hyperref}
\newcommand{\ie}{i.e.\ }
\newcommand{\cf}{cf.\ }
\newcommand{\resp}{respectively\ }
\newcommand{\Subsection}{\subsection}
\providecommand{\nobreakdash}{\nobreak\hbox{-}}
\setlength{\emergencystretch}{2em}
\multlinegap0pt
\usepackage{amssymb}
\usepackage[shortlabels]{enumitem}
\usepackage{tikz}
\newtheorem{lemma}{Lemma}
\newcommand{\C}{\mathbb C}
\newcommand{\bC}{\mathbb{C}}
\newcommand{\bN}{\mathbb{N}}
\renewcommand{\geq}{\geqslant}
\title{Isolated hypersurface singularities may be \\ stably degenerate}
\author{Mark McLean, Ivan Smith}
\date{Source: arXiv:2607.03497v1 (CC BY 4.0)}
\begin{document}
\maketitle

\noindent\emph{Source and licence.} Isolated hypersurface singularities may be \\ stably degenerate. Version arXiv:2607.03497v1 (CC BY 4.0), distributed by arXiv under the Creative Commons Attribution 4.0 International licence, http://creativecommons.org/licenses/by/4.0/, as recorded in the arXiv licence metadata for this exact version and shown on the abstract page https://arxiv.org/abs/2607.03497v1. Full licence notice: \texttt{licenses/arxiv-2607.03497-CC-BY-4.0.txt}.\par\smallskip

\noindent\textit{Editorial scope.} Counted unit: the article's lemma lemmaStrataFibration (source0.tex lines 938--943). The article's complete original proof of it is delivered with it (source0.tex lines 944--980, one \texttt{proof} environment of 37 lines). No mathematics is written, repaired or re-derived: the statement and the proof are the article's own lines, in the article's own notation, and no retained line is substituted or re-flowed.  Retained uncounted as the source-local context the proof needs: lines 895--907.  Nothing was omitted from any retained span, so the package carries no omission record.  No retained line contains a citation, so the wrapper carries no bibliography and no bibliography entry.  No retained line contains a figure environment, an image-inclusion command or drawing code, so no image is delivered and the package declares no asset dependency.  Editorial formatting only, in addition: an AMS \texttt{amsart} preamble that restates the article's own declarations and macros used by the retained lines, namely the environments \texttt{lemma} on separate counters, together with the standard packages those lines need.  The retained environments print in this document as Lemma 1 in order of appearance, whereas the article prints the same items with the numbering of its own full document.  The complete native source of the article as distributed in the arXiv e-print for this exact version is bundled unchanged as \texttt{sources/204395/source0.tex}.
\begin{lemma} \label{LemmaChiComputation}
	For each $k \in \bN$,
	$\chi_{g_k}$ is isomorphic to
	\begin{equation} \label{eqnMainChiEqn}
		\begin{aligned}
			\bC^{(d-1)N + 2k} \times \Big\{ & (u_1,\cdots,u_{d-1},v_1,\cdots,v_{d-1},w_1) \in (\bC^k)^{2d-2} \times \bC^N \ : \ \\
			                                & \sum_{j=1}^{m-1} u_j v_{m-j} = 0, \ m=2,\cdots,d-1,                               \\ &
			\sum_{j=1}^{d-1} u_j v_{d-j} + G(w_1) \neq 0
			\Big\}.
		\end{aligned}
	\end{equation}
	Also $\chi'_{g_k}$ is isomorphic to  the  variety defined by the same equations except that the last equation is zero (rather than non zero) and $\chi''_{g_k}$ is also isomorphic to the same variety, but with the last equation removed completely.
\end{lemma}

\begin{lemma} \label{lemmaStrataFibration}
	For each $i= 0,\ldots,d$,
	$\chi'_{g_k,i}-\chi'_{g_k,i+1}$ is isomorphic to the total space of an affine bundle\footnote{By an \emph{affine bundle} we mean a morphism which, Zariski locally over the base, is isomorphic to a projection map $\bC^m \times U \to U$.} over $\bC^k-0$.
	Similarly for each $i = 0,\ldots,d-1$,
	$\chi''_{g_k,i}-\chi''_{g_k,i+1}$ is isomorphic to the total space of an affine bundle over $\bC^k-0$.
\end{lemma}

\begin{proof}
	Fix a value of $i$. We use the vectors $u_1,\cdots,u_k,v_1,\cdots,v_d,w_0,\cdots,w_d$ from the proof of Lemma \ref{LemmaChiComputation}.
	Let $u_{j,l}$ be the $l$th coordinate of $u_j$ for each $j \geq i$ and each $l \in \{1,\cdots,k\}$.
	Similarly define $v_{j,l}$ for each such $j,l$.
	Define
	\begin{equation}
		v_{j,\neq l} := (v_{j,1},\cdots,v_{j,l-1},v_{j,l+1},\cdots,v_{j,k})
	\end{equation}
	for each $j \geq i$ and $l \in \{1,\cdots,k\}$.




	We claim that the map
	\begin{equation}
		u_i: \chi'_{g_k,i}-\chi'_{g_k,i+1} \to \C^k-0
	\end{equation}
	is the projection map witnessing the domain as the total space of an affine bundle.
	To do this, we will construct a trivialisation of $u_i$ in the region where $u_{i,l} \neq 0$ for some fixed $l \in \{1,\cdots,k\}$.
	By Lemma \ref{LemmaChiComputation}, we have $(\chi'_{g_k,i}-\chi'_{g_k,i+1}) \cap \{u_{i,l} \neq 0\}$ is isomorphic to
	\begin{equation} \label{eqnMainChiEqn2}
		\begin{aligned}
			\bC^{(d-1)N+2k} \times & \Big\{(u_{i+1},\cdots,u_{d-1},v_1,\cdots,v_{d-1},w_1)  \in (\bC^{k})^{2(d-1)-i} \times \bC^N \ : \
			u_{i,l} \neq 0,                                                                                                                                      \\ &
			v_{m-i,l} = -\frac{1}{u_{i,l}}\Big(\sum_{\underset{l' \neq l}{l'=1}}^k u_{i,l'}v_{m-i,l'} + \sum_{j=i+1}^{m-1} u_j v_{m-j}\Big), \ m=i+1,\cdots,d-1, \\ &
			v_{d-i,l} = -\frac{1}{u_{i,l}}\Big(\sum_{\underset{l' \neq l}{l'=1}}^k u_{i,l'}v_{d-i,l'} + \sum_{j=i+1}^{d-1} u_j v_{d-j}+G(w_1)\Big) \Big\}.
		\end{aligned}
	\end{equation}
	Therefore our Zariski local trivialisation is given by:
	\begin{equation}
		\begin{aligned}
			 & (\chi'_{g_k,i}-\chi'_{g_k,i+1}) \cap \{u_{i,l} \neq 0\} \to \bC^{(k-1)(d-i) + kd + Nd} \times (\bC^{l-1} \times \bC^* \times \bC^{k-l-1}),
			\\  & (u_i,\cdots,u_d,v_1,\cdots,v_d,w_1,\cdots,w_d) \to \\ &  ((u_{i+1},\cdots,u_d,v_{1,\neq l},\cdots,v_{d-i,\neq l},v_{d-i+1},\cdots,v_d,w_1,\cdots,w_d),u_i)
		\end{aligned}
	\end{equation}
	A similar argument holds for $\chi''_{g_k,i}-\chi''_{g_k,i+1}$.
\end{proof}

\end{document}
